% =============================================================================
% MATHEMATICAL ANALYSIS 05: Electromagnetic Factor of Two
% Stand-alone derivation for inspection
% =============================================================================

\documentclass[12pt]{article}
\usepackage{amsmath,amssymb,amsthm}
\usepackage[margin=1in]{geometry}
\usepackage{hyperref}

\newcommand{\Scfac}{a}
\newcommand{\Wight}{W}
\newcommand{\Nnum}{\mathcal{N}}
\newcommand{\pd}{\partial}
\newcommand{\be}{\begin{equation}}
\newcommand{\ee}{\end{equation}}
\newcommand{\lb}{\left}
\newcommand{\rb}{\right}

\title{Analysis 05: Electromagnetic Polarization Sum and Factor of Two}
\author{Calculation Notes --- Satishchandran \& Rao}
\date{}

\begin{document}
\maketitle

\section*{Conformal invariance of Maxwell theory in 4D}

The Maxwell action in a curved spacetime is
%
\begin{equation}
  S_{\rm EM} = -\frac{1}{4}\int F_{ab}F^{ab}\sqrt{-g}\,d^{4}x \,,
\end{equation}
%
where $F_{ab} = \pd_{a}A_{b}-\pd_{b}A_{a}$.

Under a Weyl rescaling $g_{ab}\to\Omega^{2}\hat{g}_{ab}$:
\begin{itemize}
\item $g^{ab}\to\Omega^{-2}\hat{g}^{ab}$ (raised indices gain $\Omega^{-2}$ each)
\item $F^{ab} = g^{ac}g^{bd}F_{cd}\to\Omega^{-4}\hat{F}^{ab}$
\item $F_{ab}F^{ab} = F_{ab}g^{ac}g^{bd}F_{cd}\to\Omega^{-4}\hat{F}_{ab}\hat{F}^{ab}$
\item $\sqrt{-g}\to\Omega^{4}\sqrt{-\hat{g}}$
\end{itemize}
Therefore:
%
\begin{equation}
  F_{ab}F^{ab}\sqrt{-g}\;d^{4}x
  \to
  \Omega^{-4}\hat{F}_{ab}\hat{F}^{ab}\cdot\Omega^{4}\sqrt{-\hat{g}}\,d^{4}x
  = \hat{F}_{ab}\hat{F}^{ab}\sqrt{-\hat{g}}\,d^{4}x \,.
\end{equation}
%
The action is \emph{conformally invariant in $d=4$}: no $\Omega$ factors remain.
(In general $d$ dimensions, the invariance requires $\Omega^{4-d}=1$, i.e., $d=4$.)

Consequence: the Maxwell equations in the FLRW metric $g_{ab} = \Scfac^{2}(\eta)\eta_{ab}$
reduce to the flat-space Maxwell equations for the potential $A_{\mu}$ (no conformal weight
for the $A_{\mu}$ components in Lorenz gauge).

\section*{EM Wightman function in the conformal vacuum}

Since $A_{\mu}$ satisfies the flat-space wave equation in conformal coordinates:
%
\begin{equation}
  \langle\Psi_{0}|A_{\mu}^{\rm in}(x_{1})A_{\nu}^{\rm in}(x_{2})|\Psi_{0}\rangle
  =
  W_{\mu\nu}^{\rm Mink}(\tilde{x}_{1},\tilde{x}_{2}) \,.
\end{equation}
%
In Lorenz gauge in flat space (Minkowski), the photon propagator in the Feynman/Wightman
prescription is:
%
\begin{equation}
  W_{\mu\nu}^{\rm Mink}(x,x')
  =
  -\eta_{\mu\nu}\,W_{\rm scalar}^{\rm Mink}(x,x')
  =
  \frac{\eta_{\mu\nu}}{4\pi^{2}\tilde{\sigma}_{\varepsilon}(x,x')} \,,
  \label{eq:W-EM-Mink}
\end{equation}
%
where $W_{\rm scalar}^{\rm Mink} = -1/(4\pi^{2}\tilde\sigma_\varepsilon)$ is the massless
scalar Wightman function.  Note: in FLRW the physical EM propagator has \emph{no}
$1/[\Scfac(t_{1})\Scfac(t_{2})]$ factor (unlike the scalar field $\varphi$, which has
conformal weight $-1$); the covariant vector $A_\mu$ has conformal weight 0 in $d=4$.

\section*{Decoherence exponent for EM}

The charged body couples to the EM field via $j^{\mu}A_{\mu}$.  The current difference
$j_{1}^{\mu}-j_{2}^{\mu}$ in the dipole approximation is (cf. scalar case):
%
\begin{equation}
  j_{1}^{i}(\tau,\mathbf{x}) - j_{2}^{i}(\tau,\mathbf{x})
  \approx
  q\,d_{0}\,s^{i}\,\pd_{i}\delta^{(3)}(\mathbf{x}-\mathbf{X}(\tau)) \,,
\end{equation}
%
where $i=1,2,3$ are spatial indices and the time component contributes subleadingly
in the non-relativistic limit.  The decoherence exponent is
%
\begin{equation}
  \Gamma_{\rm EM}
  =
  q^{2}d_{0}^{2}
  \int\!\!\int d\tau_{1}d\tau_{2}\;
  s^{i}s^{j}\,\pd_{i}\pd_{j'}\,
  W_{00}^{(\rm spatial)} \,,
\end{equation}
%
where $W_{ij}^{(\rm spatial)}$ is the spatial part of the EM Wightman function.

For the spatial components in Lorenz gauge (at coincident spatial points):
%
\begin{equation}
  W_{ij}^{\rm Mink}(\eta_{1},\mathbf{X};\eta_{2},\mathbf{X})
  =
  -\frac{\eta_{ij}}{4\pi^{2}\tilde\sigma_\varepsilon}\bigg|_{|\Delta\mathbf{x}|=0}
  =
  \frac{\delta_{ij}}{4\pi^{2}(\eta_{1}-\eta_{2}-i\varepsilon)^{2}} \,.
\end{equation}
%
Taking the spatial derivatives $\pd_{x^i}\pd_{x'^j}$:
%
\begin{equation}
  s^{i}s^{j}\pd_{x^i}\pd_{x'^j}\,W_{ij}\big|_{x=x'}
  =
  s^{i}s^{j}\cdot\delta_{ij}\cdot\frac{2}{(\eta_1-\eta_2-i\varepsilon)^4}
  = \frac{2|s|^{2}}{(\eta_1-\eta_2-i\varepsilon)^{4}}
  = \frac{2}{(\eta_1-\eta_2-i\varepsilon)^{4}} \,.
\end{equation}
%
(using $|s|^{2}=1$ for unit separation vector).

\section*{Polarization counting and the factor of 2}

For the \emph{scalar} field, the spatial derivative gave (Analysis~02):
%
\begin{equation}
  s^{i}s^{j}\pd_{x^i}\pd_{x'^j}\,\Wight_{\rm scalar}\big|_{x=x'}
  \approx
  \frac{2/3}{(\eta_1-\eta_2-i\varepsilon)^{4}\cdot\Scfac_1\Scfac_2} \,.
\end{equation}
%
(factor $2/3$ from $\hat{s}^i\hat{s}^j\delta_{ij}\cdot(2/\tilde\sigma^2)$ with isotropic average
$\langle s^is^j\rangle = \delta^{ij}/3$).

After including the dτ Jacobians and the $1/(\Scfac_1\Scfac_2)$ from $\Wight_{\rm scalar}$
cancellation, the scalar integral gives:
%
\begin{equation}
  \Gamma_{\rm scalar}
  =
  \frac{q^{2}d_{0}^{2}}{12\pi^{2}}
  \int\!\!\int d\eta_1 d\eta_2\;\frac{1}{(\eta_1-\eta_2-i\varepsilon)^{4}} \,.
\end{equation}

For the EM field, the EM propagator has no $1/(\Scfac_1\Scfac_2)$ factor, but the
dτ Jacobians still provide two factors of $\Scfac$.  The spatial component sum is
$W_{ij}^{\rm Mink} = \delta_{ij}\cdot W_{\rm scalar}^{\rm Mink}$ (in Lorenz gauge),
which when contracted with $s^is^j\delta_{ij}$ gives $|s|^2 W_{\rm scalar}^{\rm Mink} = W_{\rm scalar}^{\rm Mink}$.

\textbf{Naive counting:} the EM propagator at coincident points gives a factor of
$\sum_{i}\delta_{ii} = 3$ (from the trace) compared to the scalar which gives $1/3$
(after isotropic angular averaging).  But we need $s^is^j\langle\delta_{ij}-k_ik_j/k^2\rangle$
for physical (transverse) photons.

\subsection*{Physical polarization sum for transverse photons}

Physical photons are transverse: only modes perpendicular to $\mathbf{k}$ contribute.
The transverse projector is $\Pi_{ij}(\mathbf{k}) = \delta_{ij} - k_ik_j/|\mathbf{k}|^2$.
Contracting with the dipole $s^is^j$:
%
\begin{align}
  s^is^j\Pi_{ij}(\mathbf{k})
  &= s^is^j\delta_{ij} - s^is^j\frac{k_ik_j}{|\mathbf{k}|^2}
  = 1 - (\hat\mathbf{k}\cdot\hat\mathbf{s})^2 \,.
\end{align}
%
Averaging over photon propagation directions $\hat\mathbf{k}$:
%
\begin{align}
  \int\frac{d^{2}\hat{k}}{4\pi}\lb[1 - (\hat\mathbf{k}\cdot\hat\mathbf{s})^{2}\rb]
  &=
  1 - \int\frac{d^{2}\hat{k}}{4\pi}\cos^{2}\theta
  = 1 - \frac{1}{3} = \frac{2}{3} \,.
\end{align}
%
For the scalar field, the analogous integral (longitudinal projection) gives $1/3$.
Therefore:
%
\begin{equation}
  \frac{\Gamma_{\rm EM}}{\Gamma_{\rm scalar}}
  =
  \frac{2/3}{1/3} = 2 \,.
  \label{eq:ratio-2}
\end{equation}

\textbf{This is the factor of 2.}  It is purely kinematic: it counts the ratio
of transverse ($2/3$) to ``scalar'' ($1/3$) angular averages, which equals the
number of physical photon polarizations (2).

\section*{Formal derivation via polarization completeness}

In $d=4$, the sum over physical photon polarizations is
$\sum_{\lambda=1,2}\epsilon^{\mu}_{(\lambda)}(k)\epsilon^{*\nu}_{(\lambda)}(k) = -\eta^{\mu\nu} + \frac{k^\mu\bar{k}^\nu + k^\nu\bar{k}^\mu}{k\cdot\bar{k}}$,
where $\bar{k}$ is an auxiliary null vector.  For the transverse spatial components:
%
\begin{equation}
  \sum_{\lambda}\epsilon^{i}_{(\lambda)}\epsilon^{*j}_{(\lambda)}
  = \delta^{ij} - \hat{k}^{i}\hat{k}^{j}
  = \Pi^{ij}(\mathbf{k}) \,.
\end{equation}
%
Contracting with $s_{i}s_{j}$ and integrating over $\hat\mathbf{k}$:
%
\begin{align}
  \int\frac{d^3\mathbf{k}}{(2\pi)^3}\frac{s^is^j\Pi_{ij}(\hat\mathbf{k})}{2|\mathbf{k}|}
  &\propto
  \int_{0}^{\infty}dk\,k^2\cdot\frac{1}{k}\cdot\frac{2}{3}
  \propto \frac{2}{3}\int dk\,k \,.
\end{align}
%
For the scalar:
%
\begin{equation}
  \int\frac{d^3\mathbf{k}}{(2\pi)^3}\frac{s^is^j\delta_{ij}/3}{2|\mathbf{k}|}
  \propto \frac{1}{3}\int dk\,k \,.
\end{equation}
%
Ratio: $(2/3)/(1/3) = 2$ \checkmark.

\section*{Explicit EM results}

Combining the ratio \eqref{eq:ratio-2} with the scalar results:
%
\begin{align}
  \Nnum_{\rm EM}^{\rm dS}
  &= 2\Nnum_{\rm scalar}^{\rm dS}
  = \frac{q^{2}d_{0}^{2}H^{3}T}{48\pi^{2}} \,. \\[4pt]
  \Nnum_{\rm EM}^{\rm PL}
  &= 2\Nnum_{\rm scalar}^{\rm PL}
  = \frac{q^{2}d_{0}^{2}p^{3}}{24\pi^{3}}
  \cdot\frac{2T_{0}^{2}+\widetilde{T}(2+\widetilde{T})}{T_{0}^{2}(1+\widetilde{T})^{2}}
  \cdot\ln(1+\widetilde{T}) \,.
\end{align}

\section*{Universality}

The factor of 2 holds for every conformally flat FLRW background with a conformal
vacuum state, because:
\begin{enumerate}
\item Conformal invariance maps the EM problem to flat space (no $\Omega$-dependent prefactors).
\item The flat-space polarization sum contributes a factor $2/3$ vs $1/3$ for scalars.
\item This ratio is purely kinematic and background-independent.
\end{enumerate}

\section*{Summary}

$\Nnum_{\rm EM} = 2\Nnum_{\rm scalar}$ universally, counting the two physical transverse
polarizations of the photon.  The derivation is:
(1) Maxwell conformal invariance in $d=4$ removes scale-factor dependence;
(2) physical photon polarizations contribute $2/3$ angular average vs $1/3$ for scalar;
(3) ratio $= 2$.

\end{document}
